% !TEX root = ../../ctfp-print.tex

\lettrine[lhang=0.17]{モ}{ノイドのいくつかの定式化を見てきた}：集合として、単一対象圏として、モノイダル圏の対象として。この単純な概念からさらに多くのものを引き出せるだろうか？

試してみよう。モノイドを集合 $m$ と一対の関数として定義する：
\begin{align*}
  \mu  & \Colon m\times{}m \to m \\
  \eta & \Colon 1 \to m
\end{align*}
ここで 1 は $\Set$ における終対象—単一要素集合—だ。最初の関数は乗法を定義する（一対の要素を取ってその積を返す）、二番目は $m$ から単位元を選択する。これらの型シグネチャを持つ二つの関数のすべての組み合わせがモノイドになるわけではない。そのためには追加の条件を課す必要がある：結合律と単位律だ。しかし今はそれを忘れて、単に「潜在的モノイド」について考えよう。一対の関数は二つの関数集合の直積の要素だ。これらの集合は指数対象として表現できることを知っている：
\begin{align*}
  \mu  & \in m^{m\times{}m} \\
  \eta & \in m^1
\end{align*}
これら二つの集合の直積は：
\[m^{m\times{}m}\times{}m^1\]
高校代数（すべてのデカルト閉圏で成り立つ）を使ってこれを書き直せる：
\[m^{m\times{}m + 1}\]
$+$ 記号は $\Set$ における余積を表す。一対の関数を単一の関数—集合の要素—に置き換えた：
\[m\times{}m + 1 \to m\]
この関数集合の任意の要素が潜在的モノイドだ。

この定式化の美しさはそれが興味深い一般化につながることだ。例えば、この言語を使ってどのように群を記述するだろうか？群はすべての要素に逆元を割り当てる追加の関数が一つあるモノイドだ。後者は型 $m \to m$ の関数だ。例えば、整数は二項演算として加法、単位元として零、逆元として負号を持つ群を形成する。群を定義するためには三つの関数の三つ組から始める：
\begin{align*}
  m\times{}m \to m \\
  m \to m          \\
  1 \to m
\end{align*}
以前と同様に、これらの三つ組すべてを一つの関数集合に結合できる：
\[m\times{}m + m + 1 \to m\]
一つの二項演算子（加法）、一つの単項演算子（負号）、一つの零項演算子（恒等元—ここでは零）から始めた。それらを一つの関数に結合した。このシグネチャを持つすべての関数が潜在的な群を定義する。

このように続けられる。例えば、環を定義するためにはさらに一つの二項演算子と一つの零項演算子を追加する、というように。毎回、左辺が（終対象を含む場合もある零乗を含む）冪の和で、右辺が集合自体である関数型に行き着く。

さて、一般化で思い切ってみよう。まず、集合を対象に、関数を射に置き換えられる。n 項演算子を n 項積からの射として定義できる。これは有限積をサポートする圏が必要であることを意味する。零項演算子には終対象の存在が必要だ。したがってデカルト圏が必要だ。これらの演算子を結合するためには指数が必要なので、デカルト閉圏だ。最後に、代数的計算を完成させるために余積が必要だ。

あるいは、式を導いた方法を忘れて最終的な産物に集中することもできる。射の左辺にある積の和が自己関手を定義する。代わりに任意の自己関手 $F$ を選んだらどうだろうか？その場合は圏に何も制約を課す必要がない。得られるものは $F$-代数と呼ばれる。

$F$-代数は自己関手 $F$、対象 $a$、および射
\[F a \to a\]
からなる三つ組だ。対象は\newterm{担体}（carrier）、基底対象、またはプログラミングの文脈では担体型と呼ばれることが多い。射は評価関数または構造写像と呼ばれることが多い。関手 $F$ を式を形成するものとして、射をそれらを評価するものとして考えよう。

$F$-代数の Haskell の定義を示す：

\src{snippet01}
これは代数をその評価関数と同一視する。

モノイドの例では、問題の関手は次のものだ：

\src{snippet02}
これは $1 + a\times{}a$ の Haskell 表現だ（\hyperref[simple-algebraic-data-types]{代数的データ構造}を思い出してほしい）。

環は次の関手を使って定義される：

\src{snippet03}
これは $1 + 1 + a\times{}a + a\times{}a + a$ の Haskell 表現だ。

環の例として整数の集合がある。担体型として \code{Integer} を選んで評価関数を次のように定義できる：

\src{snippet04}
同じ関手 \code{RingF} に基づく $F$-代数はさらに多くある。例えば、多項式は環を形成し、正方行列も同様だ。

見てわかるように、関手の役割は代数の評価器を使って評価できる式を生成することだ。これまで非常に単純な式しか見てきていない。多くの場合、再帰を使って定義できるより複雑な式に興味がある。

\section{再帰}

任意の式木を生成する一つの方法は、関手定義内の変数 \code{a} を再帰で置き換えることだ。例えば、環における任意の式はこの木状データ構造によって生成される：

\src{snippet05}
元の環の評価器を再帰バージョンに置き換えられる：

\src{snippet06}
すべての整数を 1 の和として表現しなければならないため、これはまだあまり実用的ではないが、いざとなれば役立つ。

しかし、$F$-代数の言語を使って式木をどのように記述できるだろうか？関手の定義における自由型変数を、置換の結果で再帰的に置き換えるプロセスを何らかの形で形式化する必要がある。これをステップごとに行うことを想像してほしい。まず、深さ 1 の木を次のように定義する：

\src{snippet07}
\code{RingF} の定義の穴を \code{RingF a} によって生成された深さ 0 の木で埋めている。深さ 2 の木も同様に得られる：

\src{snippet08}
これは次のように書くこともできる：

\src{snippet09}
このプロセスを続けると、記号方程式を書ける：

\begin{snipv}
type RingF\textsubscript{n+1} a = RingF (RingF\textsubscript{n} a)
\end{snipv}
概念的には、このプロセスを無限に繰り返すと \code{Expr} に行き着く。\code{Expr} は \code{a} に依存しないことに注目してほしい。旅の出発点は問題ではなく、常に同じ場所に行き着く。これは任意の圏における任意の自己関手に対して常に成り立つわけではないが、圏 $\Set$ では都合よく成り立つ。

もちろん、これは大ざっぱな議論であり、後でより厳密にする。

自己関手を無限回適用すると\newterm{不動点}（fixed point）が生成される。それは次のように定義される対象だ：
\[\mathit{Fix}\ f = f\ (\mathit{Fix}\ f)\]
この定義の直観は、$\mathit{Fix}\ f$ を得るために $f$ を無限回適用したので、もう一回適用しても何も変わらないということだ。Haskell では、不動点の定義は次のようになる：

\src{snippet10}
おそらく、コンストラクタの名前が定義される型の名前と異なれば読みやすくなる：

\src{snippet11}
しかし、受け入れられている記法に従う。コンストラクタ \code{Fix}（または好みによっては \code{In}）は関数として見ることができる：

\src{snippet12}
関手適用の一層を剥がす関数もある：

\src{snippet13}
二つの関数は互いに逆関数だ。これらの関数は後で使う。

\section{$F$-代数の圏}

これは最古のトリックだ：新しい対象を構成する方法を思いついたら、それらが圏を形成するかどうか確認せよ。驚くことではないが、与えられた自己関手 $F$ 上の代数は圏を形成する。その圏の対象は代数—担体対象 $a$ と射 $F a \to a$ の対—であり、いずれも元の圏 $\cat{C}$ のものだ。

全体像を完成させるためには、$F$-代数の圏における射を定義しなければならない。射は一つの代数 $(a, f)$ を別の代数 $(b, g)$ に写さなければならない。それを担体を写す射 $m$ として定義する—元の圏で $a$ から $b$ へ向かう。どんな射でも良いわけではない：二つの評価器と互換性があることが必要だ。（そのような構造を保存する射を\newterm{準同型}（homomorphism）と呼ぶ。）$F$-代数の準同型を定義する方法を示す。まず、$m$ を次の写像に持ち上げられることに注目する：
\[F m \Colon F a \to F b\]
次に $g$ を続けて $b$ に到達できる。同等に、$f$ を使って $F a$ から $a$ へ向かい、次に $m$ を続けることもできる。二つのパスが等しくなることを望む：
\[g \circ F m = m \circ f\]

\begin{figure}[H]
  \centering
  \includegraphics[width=0.3\textwidth]{images/alg.png}
\end{figure}

\noindent
これが確かに圏であることを自分で納得させるのは簡単だ（ヒント：$\cat{C}$ からの恒等射はそのまま機能し、準同型の合成は準同型だ）。

$F$-代数の圏における始対象は、存在すれば\newterm{始代数}（initial algebra）と呼ばれる。この始代数の担体を $i$、評価器を $j \Colon F i \to i$ と呼ぼう。始代数の評価器 $j$ が同型射であることがわかる。この結果はランベクの定理（Lambek's theorem）として知られている。証明は始対象の定義に依存しており、それは任意の他の $F$-代数への一意の準同型 $m$ が存在することを要求する。$m$ が準同型なので、次の図式が可換でなければならない：

\begin{figure}[H]
  \centering
  \includegraphics[width=0.3\textwidth]{images/alg2.png}
\end{figure}

\noindent
次に、担体が $F i$ である代数を構成しよう。そのような代数の評価器は $F (F i)$ から $F i$ への射でなければならない。$j$ を持ち上げるだけで、そのような評価器を簡単に構成できる：
\[F j \Colon F (F i) \to F i\]
$(i, j)$ が始代数なので、そこから $(F i, F j)$ への一意の準同型 $m$ が存在しなければならない。次の図式が可換でなければならない：

\begin{figure}[H]
  \centering
  \includegraphics[width=0.3\textwidth]{images/alg3a.png}
\end{figure}

\noindent
しかし、自明に可換な図式もある（両方のパスが同じだ！）：

\begin{figure}[H]
  \centering
  \includegraphics[width=0.3\textwidth]{images/alg3.png}
\end{figure}

\noindent
これは $j$ が代数の準同型であることを示すものとして解釈でき、$(F i, F j)$ を $(i, j)$ に写す。これら二つの図式を組み合わせると次になる：

\begin{figure}[H]
  \centering
  \includegraphics[width=0.6\textwidth]{images/alg4.png}
\end{figure}

\noindent
この図式は $j \circ m$ が代数の準同型であることを示すものとして解釈できる。ただしこの場合は二つの代数が同じだ。さらに、$(i, j)$ が始代数なので、それ自体への準同型は一つしか存在できず、それは恒等射 $\id_i$ だ—それが代数の準同型であることは知っている。したがって $j \circ m = \id_i$ となる。この事実と左の図式の可換性を使って $m \circ j = \id_{Fi}$ を証明できる。これは $m$ が $j$ の逆射であることを示し、したがって $j$ は $F i$ と $i$ の間の同型射だ：
\[F i \cong i\]
しかしこれは $i$ が $F$ の不動点であることを言っているだけだ。これが最初の大ざっぱな議論の背後にある形式的な証明だ。

Haskell に戻ると：$i$ を \code{Fix f}、$j$ をコンストラクタ \code{Fix}、その逆を \code{unFix} として認識する。ランベクの定理の同型性は、始代数を得るためには関手 $f$ を取ってその引数 $a$ を \code{Fix f} で置き換えることを示している。また、不動点が $a$ に依存しない理由もわかる。

\section{自然数}

自然数も $F$-代数として定義できる。出発点は射の対だ：
\begin{align*}
  zero & \Colon 1 \to N \\
  succ & \Colon N \to N
\end{align*}
最初の射は零を選択し、二番目の射はすべての数をその後継者に写す。以前と同様に、二つを一つに結合できる：
\[1 + N \to N\]
左辺は関手を定義し、Haskell では次のように書ける：

\src{snippet14}
この関手の不動点（それが生成する始代数）は Haskell で次のようにエンコードできる：

\src{snippet15}
自然数は零か別の数の後継者だ。これは自然数のペアノ表現として知られている。

\section{カタモーフィズム}

Haskell 記法を使って始対象性の条件を書き直そう。始代数を \code{Fix f} と呼ぶ。その評価器はコンストラクタ \code{Fix} だ。始代数から同じ関手上の任意の他の代数への一意の射 \code{m} が存在する。担体が \code{a} で評価器が \code{alg} である代数を選ぼう。

\begin{figure}[H]
  \centering
  \includegraphics[width=0.4\textwidth]{images/alg5.png}
\end{figure}

\noindent
ところで、\code{m} が何であるかに注目してほしい：それは不動点の評価器、再帰式木全体の評価器だ。それを実装する一般的な方法を見つけよう。

ランベクの定理はコンストラクタ \code{Fix} が同型射であることを示す。その逆を \code{unFix} と呼んだ。したがって、この図式の矢印を一つ逆転して次を得られる：

\begin{figure}[H]
  \centering
  \includegraphics[width=0.4\textwidth]{images/alg6.png}
\end{figure}

\noindent
この図式の可換条件を書き下ろそう：

\begin{snip}{haskell}
m = alg . fmap m . unFix
\end{snip}
この方程式は \code{m} の再帰的な定義として解釈できる。再帰は関手 \code{f} を使って生成された任意の有限木に対して必ず停止する。\code{fmap m} が関手 \code{f} の最上層の下で動作することに注目することでそれがわかる。言い換えると、元の木の子ノードで動作する。子ノードは常に元の木より一段浅い。

\code{Fix\ f} を使って構成された木に \code{m} を適用すると何が起こるかを示す。\code{unFix} の作用はコンストラクタを剥がし、木の最上層を露出させる。次に、最上ノードのすべての子ノードに \code{m} を適用する。これにより型 \code{a} の結果が生成される。最後に、非再帰的な評価器 \code{alg} を適用することでそれらの結果を結合する。重要な点は、評価器 \code{alg} が単純な非再帰的関数であることだ。

任意の代数 \code{alg} に対してこれができるので、代数をパラメータとして受け取り、\code{m} と呼んだ関数を返す高階関数を定義することが理にかなっている。この高階関数は\newterm{カタモーフィズム}（catamorphism）と呼ばれる：

\src{snippet16}
その例を見てみよう。自然数を定義する関手を取る：

\src{snippet17}
担体型として \code{(Int, Int)} を選び、代数を次のように定義しよう：

\src{snippet18}
この代数のカタモーフィズム \code{cata fib} がフィボナッチ数を計算することは簡単に確認できる。

一般に、\code{NatF} に対する代数は漸化式を定義する：前の要素に関する現在の要素の値だ。カタモーフィズムはその数列の n 番目の要素を評価する。

\section{折り畳み}

\code{e} のリストは次の関手の始代数だ：

\src{snippet19}
実際、変数 \code{a} を再帰の結果（\code{List e} と呼ぶ）で置き換えると次になる：

\src{snippet20}
リスト関手に対する代数は特定の担体型を選び、二つのコンストラクタに対してパターンマッチングを行う関数を定義する。\code{NilF} に対する値は空リストをどのように評価するかを示し、\code{ConsF} に対する値は現在の要素を以前に蓄積された値とどのように結合するかを示す。

例えば、リストの長さを計算するために使える代数を示す（担体型は \code{Int}）：

\src{snippet21}
実際、結果のカタモーフィズム \code{cata lenAlg} はリストの長さを計算する。評価器はリスト要素と累積器を取って新しい累積器を返す関数と、開始値（ここでは零）の組み合わせであることに注意してほしい。値の型と累積器の型は担体型によって与えられる。

これを伝統的な Haskell の定義と比較しよう：

\src{snippet22}
\code{foldr} の二つの引数はちょうど代数の二つの成分だ。

別の例を試してみよう：

\src{snippet23}
再び、これを次と比較しよう：

\src{snippet24}
見てわかるように、\code{foldr} はカタモーフィズムのリストへの便利な特殊化にすぎない。

\section{余代数}

例によって、射の方向が逆になる F-余代数の双対構成がある：
\[a \to F a\]
与えられた関手に対する余代数も、余代数的構造を保存する準同型を持つ圏を形成する。その圏における終対象 $(t, u)$ は\newterm{終余代数}（terminal coalgebra）（または最終余代数）と呼ばれる。任意の他の代数 $(a, f)$ に対して、次の図式を可換にする一意の準同型 $m$ が存在する：

\begin{figure}[H]
  \centering
  \includegraphics[width=0.4\textwidth]{images/alg7.png}
\end{figure}

\noindent
終余代数は、射 $u \Colon t \to F t$ が同型射であるという意味で関手の不動点だ（余代数に対するランベクの定理）：
\[F t \cong t\]
終余代数はプログラミングでは通常、（無限の可能性がある）データ構造や遷移系を生成するためのレシピとして解釈される。

カタモーフィズムが始代数を評価するために使えるのと同様に、\newterm{アナモーフィズム}（anamorphism）が終余代数を余評価するために使える：

\src{snippet25}
余代数の典型的な例は、不動点が型 \code{e} の要素の無限ストリームである関手に基づいている。この関手はこれだ：

\src{snippet26}
そしてその不動点はこれだ：

\src{snippet27}
\code{StreamF e} に対する余代数は、型 \code{a} の種を取って要素と次の種からなる対（\code{StreamF} は対のおしゃれな名前だ）を生成する関数だ。

平方数のリストや逆数のような無限数列を生成する余代数の単純な例を簡単に生成できる。

より興味深い例は、素数のリストを生成する余代数だ。トリックは担体として無限リストを使うことだ。開始種はリスト \code{{[}2..{]}} になる。次の種はこのリストの尾部から 2 のすべての倍数を削除したものになる。それは 3 から始まる奇数のリストだ。次のステップでは、このリストの尾部を取って 3 のすべての倍数を削除する、というように続く。エラトステネスの篩の作りを認識できるだろう。この余代数は次の関数によって実装される：

\src{snippet28}
この余代数のアナモーフィズムは素数のリストを生成する：

\src{snippet29}
ストリームは無限リストなので、Haskell のリストに変換できるはずだ。そのために、同じ関手 \code{StreamF} を使って代数を形成し、その上でカタモーフィズムを実行できる。例えば、ストリームをリストに変換するカタモーフィズムはこれだ：

\src{snippet30}
ここで、同じ不動点が同じ自己関手に対して同時に始代数と終余代数になっている。任意の圏ではこれが常に成り立つわけではない。一般に、自己関手は多くの（またはゼロの）不動点を持つことがある。始代数はいわゆる最小不動点であり、終余代数は最大不動点だ。しかし Haskell では、どちらも同じ式で定義され、一致する。

リストのアナモーフィズムは unfold と呼ばれる。有限リストを生成するためには、関手を \code{Maybe} 対を生成するように変更する：

\src{snippet31}
\code{Nothing} の値はリストの生成を終了させる。

余代数の興味深い場合はレンズに関連している。レンズはゲッターとセッターの対として表現できる：

\src{snippet32}
ここで \code{a} は通常、型 \code{s} のフィールドを持つ何らかの積データ型だ。ゲッターはそのフィールドの値を取得し、セッターはこのフィールドを新しい値で置き換える。これら二つの関数を一つに結合できる：

\src{snippet33}
この関数をさらに書き直せる：

\src{snippet34}
ここで関手を定義した：

\src{snippet35}
これは積の和から構成された単純な代数的関手ではないことに注目してほしい。指数 $a^s$ を含む。

レンズは担体型 \code{a} を持つこの関手に対する余代数だ。\code{Store s} もコモナドであることを以前に見た。非常に行儀の良いレンズはコモナド構造と互換性がある余代数に対応することがわかる。これについては次のセクションで話す。

\section{練習問題}

\begin{enumerate}
  \tightlist
  \item
        一変数多項式の環に対する評価関数を実装せよ。多項式は $x$ の冪の前にある係数のリストとして表現できる。例えば、$4x^2-1$ は（零乗から始めて）\code{{[}-1, 0, 4{]}} として表現される。
  \item
        前の構成を $x^2y-3y^3z$ のような多変数多項式に一般化せよ。
  \item
        $2\times{}2$ 行列の環に対する代数を実装せよ。
  \item
        アナモーフィズムが自然数の平方のリストを生成する余代数を定義せよ。
  \item
        \code{unfoldr} を使って最初の $n$ 個の素数のリストを生成せよ。
\end{enumerate}
